\section{Inference-Time Reasoning: cambiar el grafo de cómputo sin modificar los pesos}

Una vez que se tiene un pretrained model, el siguiente paso no tiene por qué ser actualizar los parámetros de inmediato. Prompt, search, decomposition, code execution y note insertion pueden modificar el proceso de resolución del problema en inference time. El curso ubica estos métodos sobre un mismo eje: todos dan al modelo más cómputo intermedio o estructura externa, pero difieren en costo, verificabilidad y patrón de propagación de errores.


\lecturefigure{slide-063.jpg}{Post-training o inference-time compute: distintos puntos de intervención}{CS25 V5 official deck, p.~63}

\teachervoice{00:25:09--00:26:20, Chelsea pone en el mismo plano fine-tuning, feedback, prompting y retrieval, y señala que la CoT ofrece una ventana observable. Las notas advierten: un rationale legible no equivale necesariamente al proceso de cómputo causal real del modelo.}

\subsection{Linear Chain-of-Thought y ejemplos}

La subsección anterior dividió el post-training en varias interfaces; aquí se estudia primero la más ligera: sin modificar parámetros, solo se pide al modelo que produzca pasos intermedios. Al leer los ejemplos de CoT hay que revisar por separado la problem decomposition, la validez local de cada paso y la final answer; no se debe concluir que el razonamiento es más profundo solo porque el texto es más largo. También hay que registrar el token budget, porque una direct answer y una chain larga no implican la misma cantidad de cómputo.

La CoT sustituye el mapeo directo $x\rightarrow y$ por $x\rightarrow z_{1:k}\rightarrow y$, donde $z_i$ es un paso textual intermedio. Su beneficio puede provenir de descomponer una función compleja en transformaciones locales breves, de aumentar el test-time compute o de facilitar la localización de errores; pero los pasos intermedios también pueden ser fluidos y, aun así, erróneos.


\lecturefigure{slide-064.jpg}{Chain-of-Thought: desplegar un razonamiento complejo en pasos intermedios visibles}{CS25 V5 official deck, p.~64}


\lecturefigure{slide-065.jpg}{Ejemplo de CoT: contraste entre la respuesta directa y el razonamiento paso a paso}{CS25 V5 official deck, p.~65}

\begin{warningbox}{Faithfulness y usefulness deben evaluarse por separado}
La CoT puede mejorar la exactitud o la facilidad de depuración, pero no garantiza que cada oración refleje fielmente el mecanismo interno. La evaluación debe considerar al mismo tiempo, como mínimo, final accuracy, step validity, self-consistency y la sensibilidad a rationales irrelevantes.
\end{warningbox}

\subsection{Cómo mejorar la chain: búsqueda, supervisión y modelos pequeños}

Entre los métodos de mejora están pedir al modelo que revise lo anterior, buscar varias rutas, puntuar con un verifier, entrenar con step-level supervision o destilar reasoning traces de un modelo más fuerte. En los modelos pequeños, una chain larga puede rebasar su capacidad, por lo que la compact decomposition y la semantic understanding cobran más importancia.


\lecturefigure{slide-066.jpg}{Improving CoT: verifier, self-consistency y supervisión del proceso}{CS25 V5 official deck, p.~66}


\lecturefigure{slide-067.jpg}{Smaller models: indicaciones paso a paso, comprensión semántica y límites de capacidad}{CS25 V5 official deck, p.~67}


\lecturefigure{slide-068.jpg}{Generalizing CoT: transferencia entre dominios, entre niveles de dificultad y de estructura}{CS25 V5 official deck, p.~68}

\begin{knowledgebox}{Lectura a fondo: la CoT generalization tiene al menos cuatro sentidos}
El primero corresponde a problemas nuevos de la misma distribución, en los que el modelo reutiliza la misma plantilla de solución; el segundo, a la extrapolación de dificultad, cuando el problema requiere más pasos; el tercero, a la transferencia de dominio, de chains matemáticas a sentido común, planeación o código; el cuarto, a la transferencia estructural, que abstrae estrategias componibles a partir del formato superficial de los ejemplos. Los métodos de las slides 66--68 pueden mejorar uno de estos tipos y perjudicar otro. Al evaluar conviene variar la redacción, el orden y la información irrelevante para probar si el modelo depende del template; usar un adversarial intermediate step para comprobar si realmente lee la chain; comparar small/large models para determinar si el beneficio está limitado por la capacidad, y, por último, medir la transfer en tareas sin ejemplos. Lo que se llama «generalized CoT» solo tiene un sentido claro una vez que se explicitan estos límites.
\end{knowledgebox}

\teachervoice{00:26:25--00:29:12, en clase la CoT se amplía a un conjunto de inference procedures, no a un único prompt. Experiencia práctica: al comparar métodos hay que fijar el token budget o el wall-clock; de lo contrario, «razona mejor» puede significar solo «usó más muestreo».}

\subsection{Tree, Program y decomposition}

Tree of Thoughts genera varios candidatos en cada paso y busca entre ellos:
\[
s_{t+1}^{(j)}=f_{\theta}(s_t,a_t^{(j)}),\qquad
j^{\star}=\arg\max_j V(s_{t+1}^{(j)}).
\]
$s_t$ es el estado actual, $a_t^{(j)}$ es un thought candidato y $V$ es el evaluador. Este método convierte la generación en search, pero su desempeño depende de la proposal diversity y de la evaluator calibration.


\lecturefigure{slide-069.jpg}{Tree-of-Thoughts: generación en varias ramas, evaluación y retroceso}{CS25 V5 official deck, p.~69}


\lecturefigure{slide-070.jpg}{Program-of-Thought: el intérprete de código asume el cálculo exacto}{CS25 V5 official deck, p.~70}

Socratic decomposition, least-to-most y computation graph descomponen el problema en subproblemas dependientes; la diferencia está en si la descomposición es dialógica, secuencial o una estructura de grafo explícita. Program-of-Thought delega la aritmética al interpreter, lo que reduce la carga del modelo de lenguaje en la ejecución exacta, pero introduce la necesidad de un sandbox y de manejar errores de código.


\lecturefigure{slide-071.jpg}{Socratic questioning: indagación recursiva de los subproblemas clave}{CS25 V5 official deck, p.~71}


\lecturefigure{slide-072.jpg}{Least-to-Most prompting: resolver primero los subproblemas sencillos y luego combinarlos}{CS25 V5 official deck, p.~72}


\lecturefigure{slide-073.jpg}{Program of Thoughts: división del trabajo entre la planeación en lenguaje y el cómputo ejecutable}{CS25 V5 official deck, p.~73}


\lecturefigure{slide-074.jpg}{Computation Graphs: explicitar las dependencias como un grafo}{CS25 V5 official deck, p.~74}


\lecturefigure{slide-075.jpg}{Self-Notes: insertar en la secuencia de entrada anotaciones intermedias recuperables}{CS25 V5 official deck, p.~75}

\begin{knowledgebox}{Lectura a fondo: la aceptación de un reasoning method no puede basarse solo en la final accuracy}
En la CoT lineal hay que revisar por muestreo si cada paso es válido y si algún error se compensa por azar más adelante; en tree search hay que medir proposal diversity, branching factor, evaluator accuracy y el costo de la búsqueda; en Program-of-Thought hay que separar el planning error del execution error y restringir en el sandbox el acceso a archivos y a la red, así como el tiempo de ejecución; en decomposition hay que verificar que los subproblemas cubran el objetivo original y que el orden de dependencias esté completo; en self-notes hay que observar si los tokens posteriores realmente usan el contenido insertado. También conviene comparar con un baseline bajo el mismo token/tool budget total, porque el solo hecho de muestrear más eleva la pass rate. El informe final debe incluir como mínimo accuracy, cost, latency, verification rate, abstention y failure taxonomy. Solo así puede juzgarse si el método añade cómputo verificable o únicamente una generación más larga y más costosa.
\end{knowledgebox}


\lecturefigure{slide-076.jpg}{Latent-space CoT: añadir cómputo intermedio en el espacio oculto}{CS25 V5 official deck, p.~76}

\begin{knowledgebox}{Lectura de la figura: los dos ejes de la reasoning taxonomy}
En el eje horizontal, según el \textbf{espacio de representación}, se distinguen la chain textual, el programa, el grafo y el latent state; en el eje vertical, según la \textbf{estrategia de control}, se distinguen la generación por una sola ruta, la votación entre múltiples muestras, la búsqueda en árbol y el verifier externo. La slide 76 lleva el pensamiento al latent space, con lo que reduce el ancho de banda textual, pero sacrifica la posibilidad de inspección. Las slides anteriores, en cambio, recurren a estructuras más explícitas a cambio de capacidad de verificación. Al comparar métodos hay que fijar total sampled tokens, tool calls y evaluator compute; de lo contrario, el beneficio de los métodos de búsqueda no puede separarse de la capacidad del modelo base.
\end{knowledgebox}

\begin{table}[H]
\centering
\small
\begin{tabularx}{\textwidth}{L{0.18\textwidth} X X X}
\toprule
Método & Recurso añadido & Ventaja & Principal punto de falla \\
\midrule
Linear CoT & tokens generados adicionales & sencillo, legible & cadena de errores, alucinaciones con exceso de confianza \\
Tree search & múltiples muestras y evaluator & permite retroceder, varias rutas & costo alto, sesgo del evaluador \\
Program/Tool & ejecutor externo & cálculo exacto, verificable & seguridad, errores de interfaz y de código \\
Decomposition & estructura de subproblemas & reduce la dificultad local & descomposición errónea, dependencias omitidas \\
Latent reasoning & cómputo en estados ocultos & sin el límite del ancho de banda textual & difícil de interpretar y de supervisar \\
\bottomrule
\end{tabularx}
\caption{Clasificación de los métodos de inference-time reasoning según los recursos y los modos de falla.}
\end{table}

\begin{lstlisting}[caption={Pseudocódigo de reasoning search con presupuesto limitado}]
frontier = [initial_state(problem)]
for depth in range(max_depth):
    candidates = expand(frontier, samples_per_state)
    checked = verify_with_tools(candidates)
    frontier = top_k(checked, score="validity_then_reward")
    if solved(frontier[0]):
        return frontier[0]
return best_with_uncertainty(frontier)
\end{lstlisting}

\subsection{Resumen de la sección}

Inference-time reasoning modifica el grafo de cómputo de la resolución sin necesidad de actualizar los pretrained weights. CoT, tree, program, decomposition y latent chain añaden, respectivamente, tokens, search, herramientas, estructura o cómputo en estados ocultos; una comparación justa debe controlar el presupuesto y reportar el verifier, las herramientas y la recuperación ante fallas.
